\documentclass[11pt]{article}
\usepackage{amsmath,amssymb,amsthm,mathtools}
\usepackage[a4paper,margin=1in]{geometry}
\usepackage{enumitem}
\usepackage{hyperref}

\title{A Concrete Introduction to Cohomology on \(S^2\) and \(T^2\),\\
	with Winding Forms, Area Forms, Euler Class, and Gauss--Bonnet}
\author{}
\date{}

\newtheorem{theorem}{Theorem}
\newtheorem{proposition}{Proposition}
\newtheorem{remark}{Remark}
\newtheorem{exercise}{Exercise}

\begin{document}
	\maketitle
	
	\tableofcontents
	
	\section{What is the main idea?}
	
	This note is about one basic phenomenon:
	
	\[
	\boxed{\text{a form can be closed without being exact}.}
	\]
	
	That difference is what cohomology measures.
	
	We will study this concretely on:
	\begin{itemize}[leftmargin=2em]
		\item \(\mathbb{R}^2\setminus\{0\}\), using the winding form,
		\item \(S^2\), using the area form,
		\item \(T^2\), using the forms \(dx\), \(dy\), and \(dx\wedge dy\).
	\end{itemize}
	
	Then we will connect this to curvature and Gauss--Bonnet.
	
	The point is not to begin with abstraction. The point is to \emph{compute first}, and then see that cohomology is exactly what those computations are detecting.
	
	\section{What is de Rham cohomology?}
	
	Let \(M\) be a smooth manifold.
	
	The \(k\)-th de Rham cohomology group is
	\[
	H^k_{\mathrm{dR}}(M)
	=
	\frac{\{\text{closed }k\text{-forms on }M\}}{\{\text{exact }k\text{-forms on }M\}}.
	\]
	
	So we should keep two words separate:
	
	\begin{itemize}[leftmargin=2em]
		\item A form \(\omega\) is \emph{closed} if \(d\omega=0\).
		\item A form \(\omega\) is \emph{exact} if \(\omega=d\eta\) for some other form \(\eta\).
	\end{itemize}
	
	Every exact form is closed, because \(d^2=0\):
	\[
	\omega=d\eta \quad\Longrightarrow\quad d\omega=d(d\eta)=0.
	\]
	
	But the converse is not always true.
	
	That failure is exactly what cohomology records.
	
	So if we find a form which is closed but not exact, then we have found a nonzero cohomology class.
	
	\section{First model: the winding form on \(\mathbb{R}^2\setminus\{0\}\)}
	
	Consider the \(1\)-form
	\[
	\eta=\frac{-y\,dx+x\,dy}{x^2+y^2}
	\qquad\text{on }\mathbb{R}^2\setminus\{0\}.
	\]
	
	This is the basic winding form.
	
	\subsection{Why this is really \(d\theta\)}
	
	Write
	\[
	x=r\cos\theta,\qquad y=r\sin\theta.
	\]
	Then
	\[
	dx=\cos\theta\,dr-r\sin\theta\,d\theta,
	\qquad
	dy=\sin\theta\,dr+r\cos\theta\,d\theta.
	\]
	
	Now compute:
	\begin{align*}
		-y\,dx+x\,dy
		&=
		-r\sin\theta(\cos\theta\,dr-r\sin\theta\,d\theta)
		+r\cos\theta(\sin\theta\,dr+r\cos\theta\,d\theta)\\
		&=
		-r\sin\theta\cos\theta\,dr+r^2\sin^2\theta\,d\theta
		+r\sin\theta\cos\theta\,dr+r^2\cos^2\theta\,d\theta\\
		&=
		r^2(\sin^2\theta+\cos^2\theta)\,d\theta\\
		&=
		r^2\,d\theta.
	\end{align*}
	Since \(x^2+y^2=r^2\), we get
	\[
	\eta=d\theta.
	\]
	
	So locally, \(\eta\) looks completely harmless. It is just the differential of the angle.
	
	\subsection{Why \(\eta\) is closed}
	
	Since locally \(\eta=d\theta\), we get
	\[
	d\eta=d(d\theta)=0.
	\]
	
	So \(\eta\) is closed.
	
	\subsection{Why \(\eta\) is not exact}
	
	Now integrate around the unit circle \(\gamma(t)=(\cos t,\sin t)\), \(0\le t\le 2\pi\).
	
	Because \(\eta=d\theta\), we get
	\[
	\int_\gamma \eta=\int_0^{2\pi} d\theta=2\pi.
	\]
	
	But if \(\eta=df\) globally, then the integral of \(\eta\) around every closed curve would be \(0\):
	\[
	\int_\gamma df=0.
	\]
	
	So \(\eta\) cannot be exact.
	
	Thus
	\[
	[\eta]\neq 0\in H^1_{\mathrm{dR}}(\mathbb{R}^2\setminus\{0\}).
	\]
	
	\subsection{What is the intuition?}
	
	The reason is simple:
	
	\begin{quote}
		there is no globally single-valued angle function on \(\mathbb{R}^2\setminus\{0\}\).
	\end{quote}
	
	Locally you can choose an angle \(\theta\). Globally you cannot choose one continuously on all of \(\mathbb{R}^2\setminus\{0\}\).
	
	That is exactly why the form is closed but not exact.
	
	\section{Now move to \(S^2\)}
	
	We now look for the sphere-version of this story.
	
	For \(\mathbb{R}^2\setminus\{0\}\), the interesting object was a \(1\)-form.
	
	For \(S^2\), the interesting object will be a \(2\)-form: the area form.
	
	\subsection{Coordinates on \(S^2\)}
	
	Let \(S^2\subset \mathbb{R}^3\) be the unit sphere. Use spherical coordinates
	\[
	X(\theta,\phi)=
	(\sin\phi\cos\theta,\sin\phi\sin\theta,\cos\phi),
	\]
	where
	\[
	0<\theta<2\pi,\qquad 0<\phi<\pi.
	\]
	
	The induced metric is
	\[
	ds^2=d\phi^2+\sin^2\phi\,d\theta^2.
	\]
	
	So the area form is
	\[
	dA=\sin\phi\,d\theta\wedge d\phi.
	\]
	
	\subsection{Why \(dA\) is closed}
	
	Since \(S^2\) is \(2\)-dimensional, every \(3\)-form is automatically zero.
	
	But \(d(dA)\) is a \(3\)-form.
	
	Therefore
	\[
	d(dA)=0.
	\]
	
	So \(dA\) is closed.
	
	\subsection{Why \(dA\) is not exact}
	
	Now compute its total integral:
	\begin{align*}
		\int_{S^2} dA
		&=
		\int_0^{2\pi}\int_0^\pi \sin\phi\,d\phi\,d\theta\\
		&=
		\int_0^{2\pi}\left[-\cos\phi\right]_{\phi=0}^{\phi=\pi}\,d\theta\\
		&=
		\int_0^{2\pi}2\,d\theta\\
		&=
		4\pi.
	\end{align*}
	
	So the sphere has total area \(4\pi\).
	
	Now if \(dA=d\alpha\) globally, then by Stokes' theorem,
	\[
	\int_{S^2} dA=\int_{S^2} d\alpha=\int_{\partial S^2}\alpha=0,
	\]
	because \(S^2\) has no boundary.
	
	But we just computed
	\[
	\int_{S^2} dA=4\pi.
	\]
	
	Contradiction.
	
	So \(dA\) is not exact.
	
	Therefore
	\[
	[dA]\neq 0\in H^2_{\mathrm{dR}}(S^2).
	\]
	
	\subsection{What is the intuition?}
	
	On the punctured plane, the obstruction lived in degree \(1\), because the hole was detected by integrating a \(1\)-form around a loop.
	
	On the sphere, there is no such \(1\)-dimensional hole. So \(H^1\) vanishes.
	
	But there is still a nontrivial top-degree class: the whole sphere itself carries area. The area form cannot be the derivative of a global \(1\)-form, because its total integral is nonzero.
	
	So the sphere has:
	\[
	H^1_{\mathrm{dR}}(S^2)=0,
	\qquad
	H^2_{\mathrm{dR}}(S^2)\cong \mathbb{R}.
	\]
	
	\section{Local primitives on \(S^2\)}
	
	This is where the sphere starts to look like the winding-form story.
	
	Cover the sphere by
	\[
	U_N=S^2\setminus\{N\},
	\qquad
	U_S=S^2\setminus\{S\},
	\]
	where \(N\) is the north pole and \(S\) is the south pole.
	
	On \(U_N\), define
	\[
	\alpha_N=(1-\cos\phi)\,d\theta.
	\]
	Then
	\[
	d\alpha_N
	=
	d(1-\cos\phi)\wedge d\theta
	=
	\sin\phi\,d\phi\wedge d\theta.
	\]
	
	On \(U_S\), define
	\[
	\alpha_S=-(1+\cos\phi)\,d\theta.
	\]
	Then
	\[
	d\alpha_S
	=
	d(-(1+\cos\phi))\wedge d\theta
	=
	\sin\phi\,d\phi\wedge d\theta.
	\]
	
	So on both charts, the area form is locally a derivative:
	\[
	dA = \pm\, d\alpha_N = \pm\, d\alpha_S
	\]
	depending on orientation convention.
	
	The important point is not the sign. The important point is:
	
	\begin{quote}
		the area form is locally exact.
	\end{quote}
	
	Now compare the two local primitives on the overlap:
	\[
	\alpha_N-\alpha_S
	=
	(1-\cos\phi)d\theta+(1+\cos\phi)d\theta
	=
	2\,d\theta.
	\]
	
	So the two local primitives differ by the winding form.
	
	This is the direct analogue of the punctured-plane story.
	
	\subsection{Key message}
	
	The area form on \(S^2\) fails to be globally exact because the local primitives differ by a nontrivial angular term on the overlap.
	
	That is the sphere-version of winding.
	
	\section{Cohomology of \(S^2\)}
	
	Now we can summarize:
	
	\[
	H^0_{\mathrm{dR}}(S^2)\cong \mathbb{R},
	\qquad
	H^1_{\mathrm{dR}}(S^2)=0,
	\qquad
	H^2_{\mathrm{dR}}(S^2)\cong \mathbb{R}.
	\]
	
	Concretely:
	\begin{itemize}[leftmargin=2em]
		\item \([1]\) generates \(H^0\),
		\item there is no nontrivial \(1\)-form class,
		\item \([dA]\) generates \(H^2\).
	\end{itemize}
	
	\section{Now move to the torus \(T^2\)}
	
	The torus behaves very differently.
	
	Unlike the sphere, the torus has genuine \(1\)-dimensional winding directions.
	
	That is why \(H^1\) is nonzero.
	
	\subsection{The torus as a quotient}
	
	Let
	\[
	T^2=\mathbb{R}^2/\mathbb{Z}^2.
	\]
	
	So points \((x,y)\) and \((x+m,y+n)\) are identified whenever \(m,n\in\mathbb{Z}\).
	
	The forms \(dx\) and \(dy\) on \(\mathbb{R}^2\) descend to globally defined \(1\)-forms on \(T^2\).
	
	\subsection{Why \(dx\) and \(dy\) are closed}
	
	This is immediate:
	\[
	d(dx)=0,\qquad d(dy)=0.
	\]
	
	So both are closed.
	
	\subsection{Why \(dx\) and \(dy\) are not exact}
	
	Consider the two basic closed curves:
	\[
	\gamma_x(t)=(t,0)\mod \mathbb{Z}^2,
	\qquad
	\gamma_y(t)=(0,t)\mod \mathbb{Z}^2,
	\qquad 0\le t\le 1.
	\]
	
	Now compute:
	\[
	\int_{\gamma_x}dx=1,\qquad \int_{\gamma_x}dy=0,
	\]
	\[
	\int_{\gamma_y}dx=0,\qquad \int_{\gamma_y}dy=1.
	\]
	
	So if \(dx\) were exact, its integral around \(\gamma_x\) would have to be \(0\), but it is \(1\).
	
	Likewise \(dy\) is not exact.
	
	Therefore
	\[
	[dx],[dy]\neq 0\in H^1_{\mathrm{dR}}(T^2).
	\]
	
	Moreover, they are linearly independent, because their periods on the two basic loops distinguish them.
	
	So
	\[
	H^1_{\mathrm{dR}}(T^2)\cong \mathbb{R}^2,
	\]
	with basis \([dx],[dy]\).
	
	\subsection{What is the intuition?}
	
	Upstairs on \(\mathbb{R}^2\), we have
	\[
	dx=d(x),
	\qquad
	dy=d(y).
	\]
	
	So upstairs they are exact.
	
	But downstairs on the torus, the functions \(x\) and \(y\) are not single-valued, because
	\[
	x\mapsto x+1,\qquad y\mapsto y+1
	\]
	under the quotient identifications.
	
	So again the pattern is:
	
	\begin{quote}
		local or upstairs potentials exist, but global downstairs potentials do not.
	\end{quote}
	
	\section{The top-degree class on \(T^2\)}
	
	Now consider the \(2\)-form
	\[
	dx\wedge dy.
	\]
	
	It is closed:
	\[
	d(dx\wedge dy)=0.
	\]
	
	Now integrate over the square fundamental domain:
	\[
	\int_{T^2} dx\wedge dy
	=
	\int_0^1\int_0^1 dx\,dy
	=
	1.
	\]
	
	So \(dx\wedge dy\) cannot be exact.
	
	Therefore
	\[
	[dx\wedge dy]\neq 0\in H^2_{\mathrm{dR}}(T^2).
	\]
	
	Thus the torus has:
	\[
	H^0_{\mathrm{dR}}(T^2)\cong \mathbb{R},
	\qquad
	H^1_{\mathrm{dR}}(T^2)\cong \mathbb{R}^2,
	\qquad
	H^2_{\mathrm{dR}}(T^2)\cong \mathbb{R}.
	\]
	
	\subsection{Ring structure}
	
	The wedge product tells us how these classes multiply:
	\[
	[dx]\wedge[dx]=0,
	\qquad
	[dy]\wedge[dy]=0,
	\]
	\[
	[dx]\wedge[dy]=[dx\wedge dy],
	\qquad
	[dy]\wedge[dx]=-[dx\wedge dy].
	\]
	
	This is very concrete:
	
	\begin{quote}
		on the torus, the top class comes from wedging the two basic winding classes.
	\end{quote}
	
	\section{A local-to-global viewpoint}
	
	We now reinterpret the same ideas through overlaps.
	
	Suppose \(\omega\) is a closed \(2\)-form on a manifold \(M\).
	
	By the Poincar\'e lemma, on each small open set \(U_i\) we can write
	\[
	\omega|_{U_i}=d\alpha_i.
	\]
	
	So a closed form always has local primitives.
	
	Now compare two such primitives on an overlap:
	\[
	d(\alpha_i-\alpha_j)=0.
	\]
	
	So their difference is a closed \(1\)-form.
	
	If the overlap is simple enough, then this difference is exact:
	\[
	\alpha_i-\alpha_j=df_{ij}.
	\]
	
	So the failure to glue the local primitives is measured by functions \(f_{ij}\) on overlaps.
	
	This is the beginning of the \v{C}ech picture.
	
	Exponentiating those functions gives transition functions
	\[
	g_{ij}=e^{if_{ij}},
	\]
	which define line bundles.
	
	So cohomology, overlap data, and line bundles are all describing the same obstruction from different angles.
	
	\section{Curvature and Gauss--Bonnet}
	
	Now we turn to geometry.
	
	Let \(M\) be an oriented Riemannian surface, with Gaussian curvature \(K\) and area form \(dA\).
	
	Then the curvature \(2\)-form is
	\[
	K\,dA.
	\]
	
	\subsection{Why \(K\,dA\) is automatically closed}
	
	Because \(M\) is \(2\)-dimensional, \(d(K\,dA)\) is a \(3\)-form. So it must vanish:
	\[
	d(K\,dA)=0.
	\]
	
	So \(K\,dA\) always defines a class in
	\[
	H^2_{\mathrm{dR}}(M).
	\]
	
	\subsection{Why this is locally exact}
	
	If we choose a positively oriented orthonormal frame locally, the Levi-Civita connection is described by a local \(1\)-form \(\omega_{12}\), and the structure equation says
	\[
	d\omega_{12}=K\,dA.
	\]
	
	So curvature is locally the derivative of a connection \(1\)-form.
	
	Again the same pattern appears:
	
	\begin{quote}
		local potentials exist, but they may fail to glue globally.
	\end{quote}
	
	That failure is the Euler class.
	
	\subsection{Gauss--Bonnet}
	
	The Gauss--Bonnet theorem says:
	\[
	\int_M K\,dA=2\pi\chi(M).
	\]
	
	This means that the cohomology class of the curvature form remembers the Euler characteristic.
	
	More precisely,
	\[
	\frac{1}{2\pi}[K\,dA]
	\]
	is the Euler class in de Rham cohomology.
	
	\section{Check Gauss--Bonnet on \(S^2\)}
	
	For the unit sphere,
	\[
	K\equiv 1.
	\]
	
	So
	\[
	K\,dA=dA.
	\]
	
	Hence
	\[
	\int_{S^2}K\,dA
	=
	\int_{S^2}dA
	=
	4\pi.
	\]
	
	Since
	\[
	\chi(S^2)=2,
	\]
	we get
	\[
	2\pi\chi(S^2)=4\pi.
	\]
	
	So
	\[
	\int_{S^2}K\,dA=2\pi\chi(S^2).
	\]
	
	This is Gauss--Bonnet on the sphere.
	
	Cohomologically, the curvature class is nontrivial:
	\[
	[K\,dA]=[dA]\neq 0.
	\]
	
	\section{Check Gauss--Bonnet on \(T^2\)}
	
	Take the flat torus with metric
	\[
	ds^2=dx^2+dy^2.
	\]
	
	Then
	\[
	K\equiv 0.
	\]
	
	So
	\[
	K\,dA=0.
	\]
	
	Therefore
	\[
	\int_{T^2}K\,dA=0.
	\]
	
	Since
	\[
	\chi(T^2)=0,
	\]
	we get
	\[
	2\pi\chi(T^2)=0.
	\]
	
	So again
	\[
	\int_{T^2}K\,dA=2\pi\chi(T^2).
	\]
	
	This is Gauss--Bonnet on the flat torus.
	
	\subsection{Important distinction}
	
	The area form \(dA\) on the torus is \emph{not} exact, because its total integral is the area of the torus.
	
	But the curvature form \(K\,dA\) \emph{is} zero for the flat torus.
	
	So Gauss--Bonnet is about
	\[
	K\,dA,
	\]
	not about \(dA\) itself.
	
	This distinction matters a lot.
	
	\section{The Euler class on \(S^2\) through stereographic coordinates}
	
	Now we explain very concretely why the tangent bundle of \(S^2\) has Euler number \(2\).
	
	\subsection{North and south stereographic charts}
	
	Let
	\[
	N=(0,0,1),\qquad S=(0,0,-1).
	\]
	
	Define stereographic projection from the north pole:
	\[
	\pi_N(X,Y,Z)=\frac{X+iY}{1-Z}.
	\]
	
	Define stereographic projection from the south pole:
	\[
	\pi_S(X,Y,Z)=\frac{X-iY}{1+Z}.
	\]
	
	Let
	\[
	z=\pi_N(p),\qquad w=\pi_S(p).
	\]
	
	Then on the overlap,
	\[
	w=\frac{1}{z}.
	\]
	
	\subsection{Transition function for the tangent bundle}
	
	The tangent-bundle transition function is the derivative of the coordinate change:
	\[
	\frac{dw}{dz}=-\frac{1}{z^2}.
	\]
	
	Now restrict to the equator, where
	\[
	z=e^{i\theta}.
	\]
	
	Then
	\[
	\frac{dw}{dz}=-e^{-2i\theta}.
	\]
	
	So as \(\theta\) goes once from \(0\) to \(2\pi\), this transition function winds twice around the circle.
	
	That is the key fact.
	
	Depending on convention, one often writes the clutching map as
	\[
	e^{2i\theta},
	\]
	but the essential geometric content is always:
	
	\[
	\boxed{\text{the tangent-bundle transition winds twice}.}
	\]
	
	Therefore the Euler number is
	\[
	e(TS^2)[S^2]=2.
	\]
	
	This matches
	\[
	\chi(S^2)=2
	\]
	and also matches
	\[
	\frac{1}{2\pi}\int_{S^2}K\,dA
	=
	\frac{1}{2\pi}\cdot 4\pi
	=
	2.
	\]
	
	\section{Big comparison}
	
	Here is the whole picture in one place.
	
	\subsection*{On \(\mathbb{R}^2\setminus\{0\}\)}
	
	The winding form
	\[
	\eta=\frac{-y\,dx+x\,dy}{x^2+y^2}
	\]
	is closed but not exact.
	
	So it gives a nontrivial class in \(H^1\).
	
	\subsection*{On \(S^2\)}
	
	The area form
	\[
	dA=\sin\phi\,d\theta\wedge d\phi
	\]
	is closed but not exact.
	
	So it gives a nontrivial class in \(H^2\).
	
	There is no nontrivial \(H^1\)-class.
	
	The curvature form satisfies
	\[
	K\,dA=dA,
	\]
	and
	\[
	\int_{S^2}K\,dA=4\pi=2\pi\chi(S^2).
	\]
	
	\subsection*{On \(T^2\)}
	
	The forms \(dx\) and \(dy\) are closed but not exact, so they generate \(H^1\).
	
	Their wedge
	\[
	dx\wedge dy
	\]
	is closed but not exact, so it generates \(H^2\).
	
	For the flat torus,
	\[
	K\,dA=0,
	\]
	so the curvature class vanishes, and
	\[
	\int_{T^2}K\,dA=0=2\pi\chi(T^2).
	\]
	
	\section{Final intuition}
	
	Everything in this note can be summarized in one sentence:
	
	\begin{quote}
		cohomology measures the failure of local potentials to become global potentials.
	\end{quote}
	
	For the punctured plane, that failure appears as winding of angle.
	
	For the sphere, that failure appears in the impossibility of writing the area form as one global derivative.
	
	For the torus, that failure appears because \(x\) and \(y\) are well-defined upstairs on \(\mathbb{R}^2\), but not downstairs on the quotient \(T^2\).
	
	For curvature, that failure appears because local connection \(1\)-forms exist, but they do not globally glue. The remaining global object is the curvature \(2\)-form, and Gauss--Bonnet says its total integral is the Euler characteristic.
	
	\section{Exercises}
	
	\begin{exercise}
		Show directly that
		\[
		d\!\left(\frac{-y\,dx+x\,dy}{x^2+y^2}\right)=0
		\quad\text{on }\mathbb{R}^2\setminus\{0\}.
		\]
	\end{exercise}
	
	\begin{exercise}
		Compute
		\[
		\int_{|z|=1}\frac{-y\,dx+x\,dy}{x^2+y^2}.
		\]
		Explain why this proves the winding form is not exact.
	\end{exercise}
	
	\begin{exercise}
		Verify that on the unit sphere,
		\[
		dA=\sin\phi\,d\theta\wedge d\phi.
		\]
	\end{exercise}
	
	\begin{exercise}
		Compute
		\[
		\int_{S^2}dA=4\pi.
		\]
		Explain why this proves \(dA\) is not exact.
	\end{exercise}
	
	\begin{exercise}
		On \(S^2\setminus\{N\}\) and \(S^2\setminus\{S\}\), compute
		\[
		d((1-\cos\phi)d\theta),
		\qquad
		d(-(1+\cos\phi)d\theta).
		\]
		Then compute the difference of the two local primitives on the overlap.
	\end{exercise}
	
	\begin{exercise}
		On \(T^2=\mathbb{R}^2/\mathbb{Z}^2\), compute the periods of \(dx\) and \(dy\) on the two basic loops and prove that \([dx]\) and \([dy]\) are linearly independent.
	\end{exercise}
	
	\begin{exercise}
		Compute
		\[
		\int_{T^2}dx\wedge dy.
		\]
		Explain why this proves \(dx\wedge dy\) is not exact.
	\end{exercise}
	
	\begin{exercise}
		Check Gauss--Bonnet directly for the unit sphere and for the flat torus.
	\end{exercise}
	
	\begin{exercise}
		In the stereographic-coordinate computation, explain why the transition function has winding number \(2\) along the equator.
	\end{exercise}
	
\end{document}